\section{Part I --- Foreclosure}\label{sec:foreclosure}

\subsection{Setup}\label{sec:foreclosure:setup}

Part~I describes how a proof attempt aimed at a target can stall once
the target, the attempt and its admissible imports have been fixed.
The target is a claim \(x\in\mathsf{Claim}\) in the sense of
\Cref{sec:apparatus:closure-formation}, with interpretation
\(\llbracket x\rrbracket\); the accepted imports form a sound set
\(\mathcal B\subseteq\mathsf{Claim}\).

A proof attempt is represented by a \emph{carrier}: a presentation
through which the attempt tries to reach \(x\), together with the
residual obstruction it leaves.  The carrier also records a few
target-relative facts about itself, such as whether an accepted bridge
to the target exists.  These facts are data of the carrier, not
consequences derived here; the results of this part show what follows
from them and from explicitly stated schemas.  They are not a claim
that every possible argument for \(x\) has this form.

\begin{definition}[Typed primary carrier]\label{def:typed-primary-carrier}
A \emph{typed primary carrier} \(C\) consists of
\begin{enumerate}
\item rational matrices: an audit energy \(A_{C}\), a native-probe
  matrix \(L_{C}\), a target (dissolving) readout \(D_{C}\) and a
  declared matrix \(H_{C}\) standing for \(K_{LL}^{\dagger}\), which
  determine the carrier residual
  \[
    R_{C}=K_{DD}-K_{DL}H_{C}K_{LD},\qquad K_{FG}=FA_{C}^{\natural}G^{\mathsf T},
  \]
  where \(A_{C}^{\natural}\) inverts the nonzero diagonal entries of
  \(A_{C}\) (this is the residual of \Cref{def:carrier-residual} when
  \(A_{C}\) is diagonal and \(H_{C}=K_{LL}^{\dagger}\));
\item the verdict of its same-level audit
  (\Cref{sec:apparatus:birdint});
\item target-relative predicates on claims: \(C\leadsto x\) (aimed
  at \(x\)), \(\operatorname{NoBridge}(C,x)\) (no accepted bridge to
  the target), \(\operatorname{BridgeMismatch}(C,x)\) (structural
  incompatibility with the target), and
  \(\operatorname{MDEqTarget}(C,x)\) (the decision of its terminal
  matrix family is equivalent to the target);
\item target-independent cascade data: \(\operatorname{MRed}(C)\) (its
  cascade reduces to a family of matrix elements) and
  \(\operatorname{IRDec}(C)\) (the inherited records decide that family);
\end{enumerate}
and the proposition \(C\models\text{\TCD}\).  The \emph{residual
closure} of \(C\) is the proposition \(R_{C}=0\).
\end{definition}

The residual is used only through the proposition \(R_{C}=0\); the
matrix data serve to make it a definite proposition.  Its
interpretation as an unexplained currency in the sense of
\Cref{sec:apparatus:adequacy-residual} needs the further hypotheses
just stated, which Part~I does not use, and \(R_{C}\) is a different
object from the Hilbert-space residual \(\Xi(L)\) of
\Cref{sec:gaussian-bridge,sec:triad-bridge}.

\begin{definition}[Four typed carrier states]\label{def:four-typed-states}
Let \(x\) be a target claim and \(C\) a carrier.  \(C\) is
\emph{\Xequiv} if \((R_{C}=0)\leftrightarrow\llbracket x\rrbracket\),
and \emph{\nonXequiv} if its residual closure is equivalent to a
different claim:
\[
  \mbox{\nonXequiv}_{x}(C)
  \;\Longleftrightarrow\;
  \exists c'\in\mathsf{Claim}.\;
    \bigl((R_{C}=0)\leftrightarrow\llbracket c'\rrbracket\bigr)
    \wedge c'\neq x .
\]
The three \emph{carrier-residual-equivalent} (CRE) states, named
\emph{target-equivalent} (TE), \emph{bridge-failed} (BF) and
\emph{cascade-terminal matrix-element} (CTMT), are
\[
\begin{aligned}
  \mbox{\CRETE}_{x}(C)
    &\Longleftrightarrow
      \mbox{\Xequiv}_{x}(C)\wedge\operatorname{SelfAuditFail}(C),\\
  \mbox{\CREBF}_{x}(C)
    &\Longleftrightarrow
      \operatorname{NoBridge}(C,x)\vee\operatorname{BridgeMismatch}(C,x),\\
  \mbox{\CRECTMT}_{x}(C)
    &\Longleftrightarrow
      \operatorname{MRed}(C)\wedge\neg\operatorname{IRDec}(C)
      \wedge\neg\operatorname{MDEqTarget}(C,x):
\end{aligned}
\]
that is, the attempt can certify the target only by a failed
self-audit; it has no usable bridge to the target; or its reduction ends
at a family of matrix elements that its inherited records do not
decide.  For a claim \(c'\), put
\[
  \operatorname{NativeOutcome}_{\mathcal B}(C,c')
  \;\Longleftrightarrow\;
  \bigl(c'\in\mathcal B\wedge R_{C}=0\bigr)\vee c'\notin\mathcal B ,
\]
so that either \(c'\) is accepted and the carrier closes natively, or
\(c'\) is not accepted; we call a claim that is not accepted
\emph{open-pending}, whatever its truth value.
\end{definition}

Because claims are compared as claims, a carrier can be both
\nonXequiv\ and \Xequiv: its residual closure may be equivalent to the
target and to a different claim with the same truth value.  In a claim
language rich enough to contain, with every claim, a differently named
equivalent one, every \Xequiv\ carrier is also \nonXequiv.  The
intended reading takes claims to be statements up to a fixed
normalization, but no result below depends on this choice, and the
substance of the classification lies in the hypothesis (Cls) below
rather than in the first state.

\begin{definition}[Coverage context]\label{def:coverage-ctx}
\(\operatorname{CoverageCtx}(x,C)\) is a primitive predicate marking the
audit settings in which the four states are intended to be exhaustive.
\end{definition}

\subsection{Carrier Dichotomy}\label{sec:foreclosure:carrier-dichotomy}

Write \(\operatorname{ExactlyOne}(A,B,D)\) for
\((A\wedge\neg B\wedge\neg D)\vee(B\wedge\neg A\wedge\neg D)\vee(D\wedge\neg A\wedge\neg B)\).
The \emph{classification hypothesis} for \((x,C)\) is
\[
  \mbox{(Cls)}\qquad
  \mbox{\Xequiv}_{x}(C)\Longrightarrow
  \operatorname{ExactlyOne}\bigl(\mbox{\CRETE}_{x}(C),\mbox{\CREBF}_{x}(C),\mbox{\CRECTMT}_{x}(C)\bigr).
\]
It is not proved here.  Since the three states are defined by data
recorded in the carrier, (Cls) is a condition on carrier records, an
admissibility requirement on how attempts are recorded, rather than a conjecture about proofs.  It requires that every \Xequiv\ carrier
record exactly one failure, so a successful \Xequiv\ attempt cannot be
recorded at all; the foreclosure of this part is a consequence of this
recording convention, not evidence about the problems.  Since the
states can overlap (see above), they form, on the carriers that satisfy
the Coverage schema below, a cover rather than a partition.

\begin{theorem}[Carrier Dichotomy, conditional form]\label{thm:carrier-dichotomy}
Let \(\mathcal B\) be a sound set of accepted claims, \(x\) a target
claim, and \(C\) a carrier that satisfies \textup{(Cls)}.  Then:
\begin{enumerate}
\item if \(C\) is \nonXequiv, there is a claim \(c'\neq x\) with
  \((R_{C}=0)\leftrightarrow\llbracket c'\rrbracket\) and
  \(\operatorname{NativeOutcome}_{\mathcal B}(C,c')\): either \(c'\) is
  accepted and \(R_{C}=0\), or \(c'\) is open-pending;
\item if \(C\) is \Xequiv, exactly one of \CRETE, \CREBF\ and
  \CRECTMT\ holds.
\end{enumerate}
\end{theorem}

\begin{proof}
For (1), choose \(c'\) as in the definition.  If \(c'\in\mathcal B\),
soundness gives \(\llbracket c'\rrbracket\), hence \(R_{C}=0\); otherwise
\(c'\) is open-pending.  Part (2) is (Cls).
\end{proof}

The first branch is the only one proved from the definitions; the
second is the classification hypothesis.  The theorem needs no aiming,
typing or coverage-context hypothesis; these enter only through the
intended scope of (Cls).  The \CRETE\ alternative is
the case in which the attempt's same-level audit has failed in the
sense of \Cref{sec:apparatus:birdint}; \CREBF\ and \CRECTMT\ are the
bridge and cascade alternatives.

\subsection{Bridge Impossibility}\label{sec:foreclosure:bridge-impossibility}

A \nonXequiv\ carrier might be enriched by a bridge into a composite
that is \Xequiv.  The question is what status such a bridge has.

\begin{definition}[Carrier-side typed bridge]\label{def:carrier-side-typed-bridge}
A \emph{\TypedBridgeCarrier} is a directed transport from a carrier
toward the target side.  The primitive relation
\(\operatorname{BridgeComp}_{x}(B,C,C_{B})\) states that composing the
bridge \(B\) with the carrier \(C\) yields the composite carrier
\(C_{B}\) for the target \(x\).\footnote{This is distinct from the
calibration-side bridge of \Cref{def:calibration-side-typed-bridge}.}
\end{definition}

\begin{definition}[\(X\)-strength obligation]\label{def:x-strength-obligation}
The primitive predicate \(\operatorname{XObl}(B,x)\) states that the
bridge \(B\) carries an \emph{\XstrengthObl}: it must itself supply the
target-strength content needed for the composite to reach \(x\), and is
not merely a change of presentation.
\end{definition}

The \emph{bridge-obligation hypothesis} for \((B,C,C_{B},x)\) is
\[
\begin{aligned}
  \mbox{(BO)}\qquad
    &\operatorname{BridgeComp}_{x}(B,C,C_{B})\\
  &\quad\wedge\operatorname{ExactlyOne}\bigl(\mbox{\CRETE}_{x}(C_{B}),\mbox{\CREBF}_{x}(C_{B}),\mbox{\CRECTMT}_{x}(C_{B})\bigr)\\
  &\quad\Longrightarrow\operatorname{XObl}(B,x).
\end{aligned}
\]

\begin{corollary}[Bridge Impossibility, conditional form]\label{cor:bridge-impossibility}
Let \(\mathcal B\) be a sound set of accepted claims and \(x\) a target
claim.  Let \(C\) be a \nonXequiv\ carrier with \(R_{C}=0\), let \(B\) be
a carrier-side bridge, and let \(C_{B}\) be a composite with
\(\operatorname{BridgeComp}_{x}(B,C,C_{B})\), \(R_{C_{B}}=0\),
\((R_{C_{B}}=0)\leftrightarrow\llbracket x\rrbracket\),
\(C_{B}\leadsto x\), \(C_{B}\models\text{\TCD}\) and
\(\operatorname{CoverageCtx}(x,C_{B})\).  If \(C_{B}\) satisfies
\textup{(Cls)} and \textup{(BO)} holds, then \(\operatorname{XObl}(B,x)\).
\end{corollary}

\begin{proof}
The composite is \Xequiv, so \Cref{thm:carrier-dichotomy} applied to
\(C_{B}\) gives exactly one of its three carrier-residual-equivalent
states, and (BO) gives \(\operatorname{XObl}(B,x)\).
\end{proof}

The hypotheses of the corollary imply that \(\llbracket x\rrbracket\)
holds, since \(C_{B}\) is \Xequiv\ and \(R_{C_{B}}=0\); the corollary
describes what a bridge that actually reached the target would have to
carry.  The hypotheses are jointly satisfiable: there is a model in
which all of them hold at once, with two different claims that are both
true (\Cref{app:formalization}).  The model is degenerate (a
single carrier with zero matrices serves as both \(C\) and \(C_{B}\)),
so it shows only that the hypotheses do not contradict each other, not
that (BO) is plausible.  The proof uses only the
composition, the \(X\)-equivalence of \(C_{B}\), (Cls) and (BO); the
remaining hypotheses describe the situation the corollary is about.  The content of the corollary is
carried by (BO): a bridge that turns a \nonXequiv\ carrier into an
\Xequiv\ composite must carry the target-strength content it appears to
transport.  This is a statement inside the carrier setup of this part,
not a general impossibility result for bridges into \(x\).

\subsection{Coverage}\label{sec:foreclosure:coverage}

\begin{schema}[Coverage Theorem]\label{thm:coverage}
Let \(x\) be a target claim and \(C\) a carrier with \(C\leadsto x\),
\(C\models\text{\TCD}\) and \(\operatorname{CoverageCtx}(x,C)\).  Then
\[
  \mbox{\nonXequiv}_{x}(C)\vee
  \mbox{\CRETE}_{x}(C)\vee
  \mbox{\CREBF}_{x}(C)\vee
  \mbox{\CRECTMT}_{x}(C).
\]
\end{schema}

The schema asserts exhaustiveness only.  For a carrier that records a
claim \(r_{C}\) with \(\llbracket r_{C}\rrbracket\leftrightarrow(R_{C}=0)\),
the schema follows from (Cls): if \(C\) is not \Xequiv, then
\(r_{C}\neq x\) and \(r_{C}\) witnesses that \(C\) is \nonXequiv;
otherwise (Cls) gives one of the three states.  Without such a claim it
need not follow: in a language with a single claim, interpreted as
false, a carrier with \(R_{C}=0\) and no recorded failures satisfies
(Cls) vacuously but none of the four states.  Turning the schema into
a statement about actual proof attempts would require definitions of
the recorded predicates and a notion of a target that is open relative
to \(\mathcal B\), neither settled by accepted claims nor replaceable by
an equivalent claim that imports accepted content; this paper supplies
neither.

\subsection{Type-Finite Foreclosure}\label{sec:foreclosure:type-finite}

The set of possible carriers aimed at \(x\) may be infinite; the set of
audit outcomes is not.  Let
\(\mathsf{State}=\{\mathsf{nonX},\mathsf{TE},\mathsf{BF},\mathsf{CTMT}\}\),
and let \(\operatorname{ClassifiesAs}(s,x,C)\) mean that \(C\) is in the
state named by \(s\).

\begin{theorem}[Type-Finite Foreclosure, conditional form]\label{thm:type-finite-foreclosure}
If a carrier \(C\) satisfies the conclusion of
the Coverage schema for \(x\), then
\(\operatorname{ClassifiesAs}(s,x,C)\) for some \(s\in\mathsf{State}\).
Consequently, any family of carriers satisfying the Coverage schema for
\(x\) is classified by a map into the four-element set
\(\mathsf{State}\), however large the family is.
\end{theorem}

\begin{proof}
Choose the state named by a true disjunct of the coverage disjunction.
\end{proof}

Since the states can overlap, the classifying map is a choice and is
not canonical.

\subsection{Attack Foreclosure}\label{sec:foreclosure:attack-foreclosure-v4}

The framework has three kinds of \emph{\fim}: a \emph{\carrierPivot}
replaces the carrier by another carrier aimed at \(x\); a
\emph{\bridgeImport} composes the carrier with a bridge to \(x\); and a
\emph{\cascadeInternal} continues the reduction inside the carrier.
None of these moves is forbidden.  The claim of this subsection is that
they do not, by themselves, produce progress on \(x\) without content
from outside the framework.

\begin{definition}[Framework-scope predicate]\label{def:in-scope}
\(\operatorname{InScope}(x)\) is a primitive predicate marking the
targets to which the framework's typing and closure-formation discipline
applies.  \(\operatorname{AdvancesWithoutExternal}(m,x)\) is a
primitive predicate stating that the internal move \(m\) makes progress
on \(x\) using only the framework's internal operations.  The scope
predicate is meant to exclude targets already settled by accepted
claims, for which the schema below would fail under any natural
reading; the formal statement does not enforce this.
\end{definition}

\begin{schema}[Attack Foreclosure]\label{thm:attack-foreclosure-v4}
For every target \(x\) with \(\operatorname{InScope}(x)\) and every
framework-internal move
\(m\in\{\mbox{\carrierPivot},\mbox{\bridgeImport},\mbox{\cascadeInternal}\}\),
\[
  \neg\,\operatorname{AdvancesWithoutExternal}(m,x).
\]
\end{schema}

The intended reading is that each move either returns to the
four-state classification or exposes outside content that must be
accepted through closure formation.  This reading is an interpretation,
not a consequence of the schema: the predicates are primitive, so the
schema neither constructs a transition between carriers nor exhibits
the outside content.  The intended derivation is a case analysis.  A
pivot gives a new carrier aimed at \(x\), which
\Cref{thm:carrier-dichotomy} and the Coverage schema return to the four
states.  A bridge import that reaches an \Xequiv\ composite falls under
\Cref{cor:bridge-impossibility}.  A cascade-internal computation either
closes the residual with accepted content or enters a tower of
matrix-element terminations (\Cref{thm:ctmt-recursion}).  For a finite
chain of moves the same alternative holds move by move.  What is
missing is the theorem that this case analysis covers every admissible
chain of internal moves, which is why the statement is a schema.

The schema concerns the framework's own moves.  It is not a barrier
theorem in the sense of the classical results: relativization, by
oracles relative to which \(P=NP\) and others relative to which
\(P\neq NP\) \cite{BakerGillSolovay1975Relativization}; natural proofs,
which cannot prove superpolynomial circuit lower bounds if suitable
pseudorandom functions exist \cite{RazborovRudich1997NaturalProofs};
and algebrization, an extension of relativization to algebraic oracles
\cite{AaronsonWigderson2009Algebrization}.  The conditional results on feature-based proof
attempts in the P-versus-NP application \cite{Tsiokos2026PvNP} are
likewise conditional on explicit hypotheses and are not a barrier of
that scope.

\subsection{CTMT Recursion}\label{sec:foreclosure:ctmt-recursion}

The cascade-terminal state need not be a single step: the decision of a
terminal matrix family may itself reduce to a deeper family.  Fix a
relation \(\operatorname{Reduces}_{\sigma}(x,C,C')\), the \emph{cascade
semantics}, stating that the cascade of \(C\) reduces to that of \(C'\).
Define towers inductively: \(\mathsf{Tower}_{\sigma}(x,C,1)\) holds when
\(C\) is in \CRECTMT, and \(\mathsf{Tower}_{\sigma}(x,C,n+2)\) holds when
\(\mathsf{Tower}_{\sigma}(x,C',n+1)\) and \(\operatorname{Reduces}_{\sigma}(x,C,C')\)
for some \(C'\).

\begin{theorem}[CTMT Recursion]\label{thm:ctmt-recursion}
For every cascade semantics \(\sigma\), target \(x\) and carrier \(C\),
\[
  \mathsf{Tower}_{\sigma}(x,C,1)\Longleftrightarrow\mbox{\CRECTMT}_{x}(C),
\]
and for every \(n\ge0\), \(\mathsf{Tower}_{\sigma}(x,C',n+1)\) and
\(\operatorname{Reduces}_{\sigma}(x,C,C')\) imply
\(\mathsf{Tower}_{\sigma}(x,C,n+2)\).
\end{theorem}

\begin{proof}
A tower of depth \(1\) can only come from the first clause, since the
second produces depth at least \(2\); conversely the first clause
applies to every carrier in \CRECTMT.  The second statement is the
second clause.
\end{proof}

The theorem records that matrix-element termination can occur at any
finite depth along a chain of reductions, without introducing a fifth
state.  It does not construct a reduction in any application; the
relation \(\operatorname{Reduces}_{\sigma}\) is supplied.

\begin{table}[H]
  \centering
  \caption{The statements of Part~I and what each rests on.}
  \label{tab:foreclosure:summary}
  \small
  \renewcommand{\arraystretch}{1.1}
  \begin{tabularx}{\textwidth}{@{}>{\raggedright\arraybackslash}p{0.25\textwidth}>{\raggedright\arraybackslash}X>{\raggedright\arraybackslash}p{0.27\textwidth}@{}}
    \toprule
    Statement & Content & Status \\
    \midrule
    \Cref{thm:carrier-dichotomy} & native-outcome branch for \nonXequiv\ carriers; exactly one failure state for \Xequiv\ ones & proved, assuming (Cls) \\
    \Cref{cor:bridge-impossibility} & a bridge into an \Xequiv\ composite carries an \XstrengthObl & proved, assuming (Cls) and (BO) \\
    \Cref{thm:coverage} & the four states are exhaustive & schema; follows from (Cls) when some claim expresses \(R_{C}=0\) \\
    \Cref{thm:type-finite-foreclosure} & audit outcomes form a four-element set & proved, assuming coverage \\
    \Cref{thm:attack-foreclosure-v4} & internal moves do not advance an in-scope target without outside content & schema \\
    \Cref{thm:ctmt-recursion} & matrix-element termination nests to any finite depth & proved \\
    \bottomrule
  \end{tabularx}
\end{table}

